\begin{proposition}
\label{proposition-aic-vanishing}
Let $R$ be an absolutely integrally closed ring.
Let $M$ be a finite abelian group.
Then $H^i_\etale(\Spec(R), \underline{M}) = 0$ for $i > 0$.
\end{proposition}

\begin{proof}
Since any finite abelian group has a finite filtration whose
subquotients are cyclic of prime order, we may assume
$M = \mathbf{Z}/\ell\mathbf{Z}$ where $\ell$ is a prime number.

\medskip\noindent
Observe that all local rings of $R$ are strictly henselian, see
More on Algebra, Lemma
\ref{more-algebra-lemma-absolutely-integrally-closed-strictly-henselian}.
Furthermore, any localization of $R$ is also absolutely integrally closed
by More on Algebra, Lemma
\ref{more-algebra-lemma-absolutely-integrally-closed-quotient-localization}.
Thus Lemma \ref{lemma-gabber-criterion} tells us it suffices to
show that the kernel of
$$
H^1_\etale(D(f + g), \mathbf{Z}/\ell\mathbf{Z})
\longrightarrow
H^1_\etale(D(f(f + g)), \mathbf{Z}/\ell\mathbf{Z}) \oplus
H^1_\etale(D(g(f + g)), \mathbf{Z}/\ell\mathbf{Z})
$$
is zero for any $f, g \in R$. After replacing $R$ by $R_{f + g}$ we
reduce to the following claim: given
$\xi \in H^1_\etale(\Spec(R), \mathbf{Z}/\ell\mathbf{Z})$
and an affine open covering $\Spec(R) = U \cup V$
such that $\xi|_U$ and $\xi|_V$ are trivial, then $\xi = 0$.

\medskip\noindent
Let $A = \mathbf{Z}[\zeta]$ as above.
Since $\mathbf{Z} \subset A$ is monogenic, we can find a ring
map $A \to R$. From now on we think of $R$ as an $A$-algebra
and we think of $\Spec(R)$ as a scheme over $\Spec(A)$.
If we base change the short exact sequence (\ref{equation-ses})
to $\Spec(R)$ and take \'etale cohomology we obtain
$$
G(R) \to H(R) \to
H^1_\etale(\Spec(R), \mathbf{Z}/\ell\mathbf{Z}) \to
H^1_\etale(\Spec(R), G)
$$
Please keep this in mind during the rest of the proof.

\medskip\noindent
Let $\tau \in \Gamma(U \cap V, \mathbf{Z}/\ell\mathbf{Z})$
be a section whose boundary in the Mayer-Vietoris sequence
(Lemma \ref{lemma-mayer-vietoris}) gives $\xi$.
For $i = 0, 1, \ldots, \ell - 1$ let $A_i \subset U \cap V$
be the open and closed subset where $\tau$ has the value $i \bmod \ell$.
Thus we have a finite disjoint union decomposition
$$
U \cap V = A_0 \amalg \ldots \amalg A_{\ell - 1}
$$
such that $\tau$ is constant on each $A_i$. For
$i = 0, 1, \ldots, \ell - 1$
denote $\tau_i \in H^0(U \cap V, \mathbf{Z}/\ell\mathbf{Z})$
the element which is equal to $1$ on $A_i$ and
equal to $0$ on $A_j$ for $j \not = i$.
Then $\tau$ is a sum of multiples of the $\tau_i$\footnote{Modulo
calculation errors we have $\tau = \sum i \tau_i$.}.
Hence it suffices to show that the cohomology class corresponding
to $\tau_i$ is trivial. This reduces us to the case where
$\tau$ takes only two distinct values, namely $1$ and $0$.

\medskip\noindent
Assume $\tau$ takes only the values $1$ and $0$. Write
$$
U \cap V = A \amalg B
$$
where $A$ is the locus where $\tau = 0$ and $B$ is the locus
where $\tau = 1$. Then $A$ and $B$ are disjoint closed subsets.
Denote $\overline{A}$ and $\overline{B}$ the closures of $A$ and $B$ in
$\Spec(R)$. Then we have a ``banana'': namely
we have
$$
\overline{A} \cap \overline{B} = Z_1 \amalg Z_2
$$
with $Z_1 \subset U$ and $Z_2 \subset V$ disjoint closed subsets.
Set $T_1 = \Spec(R) \setminus V$ and $T_2 = \Spec(R) \setminus U$.
Observe that $Z_1 \subset T_1 \subset U$, $Z_2 \subset T_2 \subset V$, and
$T_1 \cap T_2 = \emptyset$. Topologically we can write
$$
\Spec(R) = \overline{A} \cup \overline{B} \cup T_1 \cup T_2
$$
We suggest drawing a picture to visualize this.
In order to prove that $\xi$ is zero, we may and do replace $R$ by
its reduction (Proposition \ref{proposition-topological-invariance}).
Below, we think of $A$, $\overline{A}$, $B$, $\overline{B}$, $T_1$, $T_2$
as reduced closed subschemes of $\Spec(R)$. Next, as scheme structures on $Z_1$
and $Z_2$ we use
$$
Z_1 = \overline{A} \cap (\overline{B} \cup T_1)
\quad\text{and}\quad
Z_2 = \overline{A} \cap (\overline{B} \cup T_2)
$$
(scheme theoretic unions and intersections as in Morphisms, Definition
\ref{morphisms-definition-scheme-theoretic-intersection-union}).

\medskip\noindent
Denote $X$ the $G$-torsor over $\Spec(R)$ corresponding to the image
of $\xi$ in $H^1(\Spec(R), G)$. If $X$ is trivial, then $\xi$
comes from an element $h \in H(R)$ (see exact sequence of cohomology
above). However, then by Lemma \ref{lemma-lift-points-H-to-G}
the element $h$ lifts to an element of $G(R)$ and we conclude
$\xi = 0$ as desired. Thus our goal is to prove that $X$ is trivial.

\medskip\noindent
Recall that the embedding $\mathbf{Z}/\ell \mathbf{Z} \to G(R)$
sends $i \bmod \ell$ to $\sigma_i \in G(R)$. Observe that $\overline{A}$
is the spectrum of an absolutely integrally closed ring (namely
a qotient of $R$). By
Lemma \ref{lemma-interpolate} we can find $g \in G(\overline{A})$ with
$g|_{\overline{A} \cap Z_1} = \sigma_0$ and
$g|_{\overline{A} \cap Z_2} = \sigma_1$ (scheme theoretically).
Then we can define
\begin{enumerate}
\item $g_1 \in G(U)$ which is $g$ on $\overline{A} \cap U$, which
is $\sigma_0$ on $\overline{B} \cap U$, and $\sigma_0$ on $T_1$, and
\item $g_2 \in G(V)$ which is $g$ on $\overline{A} \cap V$, which
is $\sigma_1$ on $\overline{B} \cap V$, and $\sigma_1$ on $T_2$.
\end{enumerate}
Namely, to find $g_1$ as in (1) we glue the section
$\sigma_0$ on $\Omega = (\overline{B} \cup T_1) \cap U$
to the restriction of the section $g$ on $\Omega' = \overline{A} \cap U$.
Note that $U = \Omega \cup \Omega'$ (scheme theoretically)
because $U$ is reduced and $\Omega \cap \Omega' = Z_1$ (scheme theoretically)
by our choice of $Z_1$. Hence by
Morphisms, Lemma \ref{morphisms-lemma-scheme-theoretic-union}
we have that $U$ is the pushout of $\Omega$ and $\Omega'$
along $Z_1$. Thus we can find $g_1$.
Similarly for the existence of $g_2$ in (2).
Then we have
$$
\tau = g_2|_{A \cup B} - g_1|_{A \cup B} \quad(\text{addition in group law})
$$
and we see that $X$ is trivial thereby finishing the proof.
\end{proof}